\subsection{LINEARITY OF DIFFERENTIATION}

PROPOSITION 1. The set of vector functions defined on an interval \(I \subset R\) , taking values in a given topological vector space E, and differentiable at the point \(x_{0}\) , is a vector space over R, and the map \(\mathbf{f} \mapsto \mathrm{Df}(x_{0})\) is a linear mapping of this space into E.

In other words, if f and g are defined on I and differentiable at the point \(x_{0}\) , then \(f + g\) and fa (a an arbitrary scalar) are differentiable at \(x_{0}\) and their derivatives there are \(\mathbf{f}'(x_{0}) + \mathbf{g}'(x_{0})\) and \(\mathbf{f}'(x_{0})a\) respectively. This follows immediately from the continuity of \(x + y\) and of xa on E × E and E respectively.

COROLLARY. The set of vector functions defined on an interval I, taking values in a given topological vector space E, and differentiable on I, is a vector space over R, and the map \(f \mapsto Df\) is a linear mapping of this space into the vector space of mappings from I into E.

Remark. If one endows the vector space of mappings from I into E and its subspace of differentiable mappings (cf.~Gen.~Top., X, p.~277) with the topology of simple convergence (or the topology of uniform convergence), the linear mapping \(\mathbf{f} \mapsto \mathrm{Df}\) is not continuous (in general) *for example, the sequence of functions \(\mathbf{f}_n(x) = \sin n^2 x / n\) converges uniformly to 0 on \(\mathbf{R}\) , but the sequence of derivatives \(\mathbf{f}_n'(x) = n\cos n^2 x\) does not converge even simply to \(0\)*

PROPOSITION 2. Let E and F be two topological vector spaces over R, and u a continuous linear map from E into F. If f is a vector function defined on an interval \(I \subset R\) , taking values in E, and differentiable at the point \(x_{0} \in I\) , then the composite function \(u \circ f\) has a derivative equal to \(\mathbf{u}(\mathbf{f}'(x_{0}))\) at \(x_{0}\) .

Indeed, since \(\frac{\mathbf{u}(\mathbf{f}(x)) - \mathbf{u}(\mathbf{f}(x_0))}{x - x_0} = \mathbf{u}\left(\frac{\mathbf{f}(x) - \mathbf{f}(x_0)}{x - x_0}\right)\) , this follows from the continuity of \(\mathbf{u}\) .

COROLLARY. If \(\varphi\) is a continuous linear form on \(E\) , then the real function \(\varphi \circ \mathbf{f}\) has a derivative equal to \(\varphi(\mathbf{f}'(x_0))\) at the point \(x_0\) .

Examples. 1) Let \(\mathbf{f} = (f_i)_{1\leqslant i\leqslant n}\) be a function with values in \(\mathbf{R}^n\) , defined on an interval \(I\subset \mathbf{R}\) ; each real function \(f_{i}\) is none other than the composite function \(\mathrm{pr}_i\circ \mathbf{f}\) , so is differentiable at the point \(x_0\) if \(\mathbf{f}\) is, and, if so, \(\mathbf{f}'(x_0) = (f_i'(x_0))_{1\leqslant i\leqslant n}\) .

*2) In Kinematics, if \(\mathbf{f}(t)\) is the position of a moving point M at time \(t\) , if \(\mathbf{g}(t)\) is the position at the same instant of the projection \(\mathbf{M}'\) of M onto a plane P (resp. a line D) with kernel a line (resp. a plane) not parallel to P (resp. D), then \(\mathbf{g}\) is the composition of the projection \(\mathbf{u}\) of \(\mathbf{R}^3\) onto P (resp. D) and of \(\mathbf{f}\) ; since \(\mathbf{u}\) is a (continuous) linear mapping one sees that the projection of the velocity of a moving point onto a plane (resp. a line) is equal to the velocity of the projection of the moving point onto the plane (resp. line).*

\begin{enumerate}
\def\labelenumi{\arabic{enumi})}
\setcounter{enumi}{2}
\tightlist
\item
  Let \(f\) be a complex-valued function defined on an interval \(\mathbf{I} \subset \mathbf{R}\) , and let \(a\) be an arbitrary complex number; prop. 2 shows that if \(f\) is differentiable at a point \(x_0\) then so is \(af\) , and the derivative of this function at \(x_0\) is equal to \(af'(x_0)\) .
\end{enumerate}

